\chapter{Getting started}\label{ch:start}
\begin{goals}
\textbf{In this chapter} you open wuffle, find your way around the desk, and draw your first data.
\end{goals}

\section{Opening wuffle}\label{sec:opening}

wuffle is the file \file{wuffle.html}. Open it in a current browser (Chrome, Edge, Firefox or Safari)
by double-clicking it or from the web address it is published at. Nothing is installed.

Which of its two homes you get depends on where it opens:

\begin{itemize}
\item On a narrow screen (a phone held upright) or inside the lead-acid Android shell, wuffle opens as the
\emph{field} compass.
\item Anywhere else it opens as the \emph{desk}.
\item A page that carries a project (\cref{sec:save-page}) always opens on the desk, into its project.
\end{itemize}

You can choose explicitly by adding \code{?field} or \code{?desk} to the address. From the desk,
\menu{View>Field mode (the compass)} switches to the compass; on a computer it takes typed measurements
instead of reading a sensor (\cref{sec:manual-entry}).

\section{The desk}\label{sec:desk}

\shot{desk-tour}{The desk with the Tocher sample open: the layers tree on the left, the net in the middle,
and the table, statistics and console tabs below.}

The desk is one window in three regions, each made of tabs you can drag and rearrange:

\begin{description}
\item[The layers tree] (left). The diagrams, each opening into its layers, and each layer into what it
draws: poles, great circles, contours, the mean, eigenvectors, labels. A \ui{data} section at the foot
lists every dataset, whether or not it is drawn. Double-click a layer for its properties.
\item[The diagrams] (middle). One tab per net, rose diagram, fabric plot or figure. The tab in front is
the \term{active diagram}: the menus and the console act on it.
\item[Table, statistics, console] (below). The rows of the dataset you clicked, the statistics of the
active diagram's layers, and the shell.
\end{description}

The menus run along the top: \menu{File}, \menu{Edit}, \menu{Diagram}, \menu{Figure}, \menu{View} and
\menu{Help}. The badge at the right says \ui{● local} (\cref{sec:local}).

\section{Your first data}\label{sec:first-data}

\begin{tryit}
\begin{enumerate}
\item Choose \menu{File>Open the tocher sample}. 200 planes open as a dataset called \code{tocher} and
are drawn as poles on \ui{Net 1}.
\item In the layers tree, open \ui{tocher} on \ui{Net 1} and tick \ui{contours}. The poles are contoured.
\item Click a pole. Its row is selected: it turns violet on the net and the table shows it.
\item Press \key{Ctrl}+\key{Z}. The contours go off again: \menu{Edit} names the step it undid.
\end{enumerate}
\end{tryit}

To open your own data, choose \menu{File>Open data or a project…} (\key{Ctrl}+\key{O}) or drop files on
the window. A delimited text file (CSV, TSV, or columns separated by spaces or semicolons) becomes a
dataset, drawn on the active net. wuffle reads the header to find the attitude:

\begin{itemize}
\item dip direction and dip, or strike and dip by the right-hand rule, as planes;
\item trend and plunge as lines;
\item a plane and a line on the same row as faults, and a direction with an angle as small circles;
\item one column holding a whole attitude, such as \code{045/30}, \code{310→70} or \code{N30E 40SE}.
\end{itemize}

wuffle never refuses a file for being odd. Rows it can't read are counted and left out of the drawing,
and when it had to guess which columns hold the attitude (no header, say), the table's heading says
how it read them.

\section{The console, briefly}

Press \key{`} (the backquote key) or choose \menu{View>Console}. The console is a shell, and its
working directory is the project itself. \cref{ch:console} explains it; for now, try this:

\begin{tryit}
\begin{enumerate}
\item With the Tocher sample open, type \code{stats} and press \key{Enter}. The count, mean, $\kappa$,
$\alpha_{95}$ and eigenvalues of the first visible layer print.
\item Type \code{ls data} to list the datasets as files, then \code{head -3 data/tocher.csv} to read the
first rows.
\end{enumerate}
\end{tryit}
